\section*{Preamble}

We conclude this journey by presenting an example of a secure integration of
\acsp{TEE} into existing protocols and applications, taking into account the
challenges presented in \cref{ch:time} and demonstrating how carefully-studied
designs can provide benefits such as increased security, efficiency and
scalability to a system. In particular, we focus on credential management
systems in \ac{V2X} applications, where credentials need to be distributed to
vehicles in a network in order to secure \ac{V2V} communication. A crucial
aspect of these systems is \emph{credential revocation}, i.e., rendering the
credentials of malicious or misbehaving vehicles unusable to preserve the safety
of the \ac{V2X} network. It is especially important that revocation is
\emph{timely}, i.e., it is executed within a short and predictable time, but
also \emph{efficient}, i.e., it scales well with the potentially large size of
the network without incurring prohibitive costs and latencies.

Existing standards on \ac{V2X} credential management raise some concerns on both
the timeliness and efficiency of their revocation
schemes~\cite{yoshizawa2022v2x_survey}. For this reason, academic proposals have
tried to integrate trusted computing technologies inside vehicles to enable a
``decentralized'' credential management and
revocation~\cite{forster2015rewire,whitefield2017privacy,desmoulins2019practical,hicks2020vehicular,larsen2021daa}.
In these schemes, each vehicle has a \emph{trusted component} responsible for
securely storing credentials and performing cryptographic operations such as
message signing and verification. Other software in the vehicle, therefore, has
no direct access to those credentials, which remain protected even in case the
vehicle becomes compromised. However, as long as credentials are valid, an
attacker is still able to request the trusted component to sign malicious
\ac{V2V} messages and potentially harm the system. To effectively make
credentials unusable, researchers have introduced the idea of
\emph{self-revocation}, where the trusted component inside the vehicle
autonomously deletes its credentials if requested by the \ac{V2X} infrastructure
via a revocation message. In these proposals, trusted components are typically
\acsp{TPM}, but \acsp{TEE} like Intel \acs{SGX} are another viable
option~\cite{jangid2022towards}.

Our contribution addresses a problem that was overlooked in previous work: What
happens if the revocation message is never received by the trusted component,
for example due to network malfunctions or malicious activity (either in the
network or inside the vehicle)? To deal with this issue, researchers proposed to
use keep-alive messages and timeouts to detect
misbehaviour~\cite{forster2015rewire,whitefield2017privacy}, but did not give
enough details to effectively solve the problem. Besides, they failed to
consider a crucial point, i.e., that trusted computing technologies, and
\acsp{TEE} in particular, cannot provide strong availability guarantees nor a
reliable notion of time (cf.~\cref{ch:time}). In addition, previous work did not
provide a comprehensive evaluation of their design, and it is unclear whether
the self-revocation approach indeed gives actual benefits compared to
traditional revocation schemes. Therefore, both the timeliness and efficiency
properties are somewhat questionable in earlier work.

The first time I read about credential management in \ac{V2X} was in January
2022, when I came across a paper from Whitefield et
al.~\cite{whitefield2017privacy}. In that occasion, I immediately recognized the
potential of self-revocation, but I also quickly realized that this approach had
many challenges that were not sufficiently addressed in the paper. Sadly, I
noticed that even subsequent work seemed to neglect these
issues~\cite{desmoulins2019practical,larsen2021daa}. This, nevertheless, turned
into an opportunity to contribute to this domain. Together with Christoph
Baumann, we designed our secure version of the self-revocation scheme, which we
refined over the months and even submitted a patent application with title
\emph{``Secure revocation mechanisms for V2X attestation protocols''} that was
filed in June 2022. The core of our work, however, was ensuring that our
approach actually gives the guarantees that earlier work could not provide. To
this extent, much time and effort was dedicated to formally verify our scheme
with Tamarin prover~\cite{meier2013tamarin} to satisfy the \emph{timeliness}
property, i.e., that we can actually guarantee timely credential revocation even
with a strong network attacker and considering the availability limitations of
\acsp{TEE}. Part of the formal verification work was also the basis to evaluate
the \emph{efficiency} of our scheme where, together with Fritz Alder and
Christoph Baumann, we demonstrated that our design is efficient and scalable.

Our evaluation is a significant contribution to validate our design and provide
enough evidence about the feasibility of our approach. For reproducibility of
results, all our code was open-sourced to allow everyone who is interested in
our work to verify our claims. Our code was also peer-reviewed as an artifact
alongside our paper and received all three ``Available'', ``Functional'', and
``Reproduced'' badges. Subsequently, we also received the ``Distinguished
Artifact Award'' during the NDSS 2024 conference. This was the recognition of
countless hours spent on improving and documenting the code to reduce as much as
possible the effort required to reproduce our results, which was particularly
appreciated by the reviewers. Our code was also presented as a demo at the
Ericsson Research Security Day 2022 and the VehicleSec 2024 workshop, co-located
with NDSS~\cite{scopelliti2024demo}.

Lastly, this contribution led to two master's theses aiming to perform a
quantitative evaluation of revocation schemes in \ac{V2X}. Indeed, due to the
lack of publicly available reference implementations, our paper could only
provide a \emph{qualitative} comparison between our self-revocation scheme and
traditional approaches. In order to carry out a ``fair'' evaluation, instead, it
is necessary that each scheme is evaluated under the same conditions to give a
meaningful comparison. To achieve this, our students implemented credential
management on top of the Artery \ac{V2X} simulation
framework~\cite{riebl2015artery}, which they used to compare active, passive and
self-revocation mechanisms. Their results show that our self-revocation scheme
gives the most balanced solution in terms of revocation time and network
overhead~\cite{dehaes2024thesis,orban2024thesis}.